\documentclass[11pt]{article}
\usepackage[a4paper,margin=24mm]{geometry}
\usepackage[T1]{fontenc}
\usepackage{lmodern,amsmath,amssymb,amsthm,mathtools,microtype,enumitem}
\usepackage[colorlinks=true,linkcolor=blue,urlcolor=blue]{hyperref}
\hypersetup{pdftitle={zeta(3) is irrational: a proof by a maximum principle},pdfauthor={Mathathon peer-review digest}}
\newcommand{\Z}{\mathbb Z}
\DeclareMathOperator{\lcm}{lcm}
\newtheorem{theorem}{Theorem}
\newtheorem{lemma}{Lemma}
\setlist{itemsep=1pt,topsep=3pt}
\title{$\zeta(3)$ is irrational: a proof by a maximum principle}
\author{Mathathon peer-review digest (Claude)}
\date{14 November 2026}
\begin{document}
\maketitle

\noindent This is a complete proof in five lemmas, using only sums of rational functions. Nothing is
integrated and no asymptotic estimate is made. The first two lemmas are classical. The way sign and
size are obtained (Lemma~\ref{lem:max}) is the point of the note; its relation to the literature is
stated at the end. Checked by AI agents only; no human has read it.

\paragraph{Plan.} Suppose $\zeta(3)=P/Q$. We build numbers $I_n$ such that (i) $\frac12Q\,D_n^3I_n$
is an integer, where $D_n=\lcm(1,\dots,n)$; (ii) $0<I_n\le2\zeta(3)\,33^{-n}$; (iii) $D_n<3^n$ for
infinitely many $n$. Then a positive integer is smaller than $Q\zeta(3)(27/33)^n$, which is absurd
for large $n$.

\paragraph{The numbers.} For $n\ge0$ let
\[
r_n(t)=\frac{(t-1)(t-2)\cdots(t-n)}{t(t+1)\cdots(t+n)},\qquad R_n(t)=r_n(t)^2,\qquad
I_n=-\sum_{t=1}^{\infty}R_n'(t).
\]
\emph{Why this function.} Its terms $\frac1{(t+k)^2}$ become $\frac1{(t+k)^3}$ after one
differentiation, which is what $\zeta(3)=\sum t^{-3}$ needs; and its double zeros at $t=1,\dots,n$
make the first $n$ terms of the series vanish. (The main document derives $r_n$ from an orthogonal
polynomial; here we simply use it.) The series converges absolutely since $R_n'(t)=O(t^{-3})$.
We write $H_k^{(j)}=\sum_{l=1}^kl^{-j}$.

\begin{lemma}[Arithmetic]\label{lem:arith}
$I_n=2(a_n\zeta(3)-b_n)$ with $a_n=\sum_{k=0}^n\binom nk^2\binom{n+k}k^2\in\Z$ and
$D_n^3b_n\in\Z$. In particular $I_0=2\zeta(3)$.
\end{lemma}
\begin{proof}
Partial fractions: $r_n(t)=\sum_{k=0}^n\frac{c_k}{t+k}$ with
$c_k=\bigl[(t+k)r_n(t)\bigr]_{t=-k}=(-1)^{n-k}\binom nk\binom{n+k}k\in\Z$. Squaring, and using
$\frac1{(t+k)(t+j)}=\frac1{j-k}\bigl(\frac1{t+k}-\frac1{t+j}\bigr)$ for $j\ne k$,
\[
R_n(t)=\sum_{k=0}^n\Bigl[\frac{c_k^2}{(t+k)^2}+\frac{e_k}{t+k}\Bigr],\qquad
e_k=2c_k\sum_{j\ne k}\frac{c_j}{j-k}.
\]
Here $D_ne_k\in\Z$, and $\sum_ke_k=0$ because the terms $(j,k)$ and $(k,j)$ cancel. Differentiating
and summing over $t\ge1$, with $\sum_{t\ge1}(t+k)^{-j}=\zeta(j)-H_k^{(j)}$:
\[
I_n=2\Bigl(\sum_kc_k^2\Bigr)\zeta(3)+\Bigl(\sum_ke_k\Bigr)\zeta(2)-\sum_k\bigl[2c_k^2H_k^{(3)}+e_kH_k^{(2)}\bigr].
\]
The middle term vanishes; $a_n=\sum c_k^2$; and $b_n=\sum_k[c_k^2H_k^{(3)}+\frac12e_kH_k^{(2)}]$
has denominator dividing $D_n^3$.
\end{proof}

\begin{lemma}[Recurrence]\label{lem:rec}
For $n\ge1$,
\begin{equation}\label{star}
(n+1)^3\bigl(I_{n+1}-I_n\bigr)-n^3\bigl(I_n-I_{n-1}\bigr)=4(2n+1)^3\,I_n .
\end{equation}
\end{lemma}
\begin{proof}
Let $\sigma_n(t)=4(2n+1)\bigl(-2t^2+t+(2n+1)^2\bigr)$ and $\beta_n=(n+1)^3+n^3+4(2n+1)^3$. We claim
\begin{equation}\label{cert}
(n+1)^3R_{n+1}(t)-\beta_nR_n(t)+n^3R_{n-1}(t)=\sigma_n(t+1)R_n(t+1)-\sigma_n(t)R_n(t).
\end{equation}
Given \eqref{cert}, differentiate and sum over $t=1,\dots,T$: the right side telescopes to
$(\sigma_nR_n)'(T+1)-(\sigma_nR_n)'(1)$. The last term is $0$ because $R_n$ has a double zero at
$t=1$; the other tends to $0$ because $\sigma_nR_n$ is a rational function with a finite limit at
infinity. So $(n+1)^3I_{n+1}-\beta_nI_n+n^3I_{n-1}=0$, which is \eqref{star}.

To prove \eqref{cert}, divide by $R_n(t)$, using
$\frac{R_{n+1}(t)}{R_n(t)}=\bigl(\frac{t-n-1}{t+n+1}\bigr)^2$,
$\frac{R_{n-1}(t)}{R_n(t)}=\bigl(\frac{t+n}{t-n}\bigr)^2$,
$\frac{R_n(t+1)}{R_n(t)}=\frac{t^4}{(t-n)^2(t+n+1)^2}$, and multiply by $(t-n)^2(t+n+1)^2$:
\begin{multline*}
(n+1)^3(t-n-1)^2(t-n)^2-\beta_n(t-n)^2(t+n+1)^2+n^3(t+n)^2(t+n+1)^2\\
=\sigma_n(t+1)\,t^4-\sigma_n(t)\,(t-n)^2(t+n+1)^2 .
\end{multline*}
Both sides are polynomials in $t$ of degree at most $4$ (on the right the terms in $t^6$ and $t^5$
cancel). So it suffices to check five zeros of the difference, counted with multiplicity. Put
$\mu=2n+1$ and $u(t)=-2t^2+t+\mu^2$, so $\sigma_n=4\mu u$; then $u(n+1)=n\mu$, $u(-n)=(n+1)\mu$,
$u'(n+1)=-(4n+3)$, $u'(-n)=4n+1$, and $\beta_n=\mu(17n^2+17n+5)$.
\begin{itemize}
\item $t=n$: both sides equal $4n^5\mu^2$. \quad $t=-n-1$: both sides equal $4(n+1)^5\mu^2$.
\item $t=0$: left $n^2(n+1)^2[(n+1)^3+n^3-\beta_n]=-4\mu^3n^2(n+1)^2$, which is the right side.
\item derivative at $t=n$: both sides equal $4n^4\mu(4n+1)$.
\quad derivative at $t=-n-1$: both sides equal $-4(n+1)^4\mu(4n+3)$.\qedhere
\end{itemize}
\end{proof}
\noindent\emph{Where \eqref{cert} comes from.} One asks for constants $A,B,C$ and a quadratic
$\sigma$ with $AR_{n+1}+BR_n+CR_{n-1}=\sigma(t+1)R_n(t+1)-\sigma(t)R_n(t)$. That is five linear
equations in six unknowns, so a solution must exist; it is unique up to scale and is the one above.
The main document (Section~10) explains this. The only new remark is that the middle coefficient,
usually written $(2n+1)(17n^2+17n+5)$, equals $(n+1)^3+n^3+4(2n+1)^3$.

\begin{lemma}[Boundedness]\label{lem:bdd}
$|I_n|<2$ for $n\ge1$.
\end{lemma}
\begin{proof}
For integers $t>n$: $0<r_n(t)<\frac1t$, since each factor $\frac{t-j}{t+j}$ is less than $1$. And
$R_n'=2R_n\varphi_n$ with $\varphi_n(t)=\sum_{j=1}^n\frac1{t-j}-\sum_{j=0}^n\frac1{t+j}$, a difference
of two positive numbers each at most $n$. As $R_n'(t)=0$ for $t\le n$,
$|I_n|\le2n\sum_{t>n}t^{-2}<2$.
\end{proof}

\paragraph{The idea of the next lemma.} Read $I_{n+1}-I_n$ as a derivative. Then \eqref{star} is the
discrete form of $(p\,u')'=q\,u$ with $p,q>0$. A solution of such an equation cannot have a positive
maximum inside its domain: there $u'=0$ and $u''\le0$, so the left side is $\le0$ while the right
side is positive. So a positive solution that ever turns upward keeps rising for ever. A solution
that stays bounded has no choice: it is positive and decreasing all the way down.

\begin{lemma}[Maximum principle]\label{lem:max}
Let $(X_n)_{n\ge0}$ be a bounded real sequence with $X_0>0$ that satisfies \eqref{star}, that is,
$(n+1)^3(X_{n+1}-X_n)-n^3(X_n-X_{n-1})=4(2n+1)^3X_n$ for $n\ge1$. Then
$X_0>X_1>X_2>\dots>0$, and $X_n<X_{n-1}/33$ for all $n\ge1$.
\end{lemma}
\begin{proof}
Put $d_n=X_{n+1}-X_n$. The equation reads $(n+1)^3d_n=n^3d_{n-1}+4(2n+1)^3X_n$ for $n\ge1$.

\emph{(a) Once up, always up.} Suppose $X_m>0$ and $d_m\ge0$ for some $m\ge0$. Then $X_{m+1}>0$, and
the equation at $n=m+1$ gives $d_{m+1}>0$. Repeating, $d_n>0$ and $X_n\ge X_{m+1}=:c>0$ for all
$n>m$. Then, as $n^3d_{n-1}\ge0$, we get $(n+1)^3d_n\ge4(2n+1)^3X_n\ge4(2n+1)^3c>4(n+1)^3c$, so
$d_n>4c$ for all $n>m$: the sequence is unbounded, a contradiction.

\emph{(b) Once non-positive, always down.} Suppose $X_m>0\ge X_{m+1}$ for some $m\ge0$. Then $d_m<0$,
and the equation at $n=m+1$ gives $(m+2)^3d_{m+1}=(m+1)^3d_m+4(2m+3)^3X_{m+1}<0$, so
$X_{m+2}<X_{m+1}\le0$. Repeating, $d_n<0$ and $X_n\le X_{m+2}=:-c<0$ for $n\ge m+2$, and then
$(n+1)^3d_n<-4(2n+1)^3c$ gives $d_n<-4c$: unbounded again, a contradiction.

\emph{(c)} If $X_m>0$ then $X_{m+1}<X_m$ by (a) and $X_{m+1}>0$ by (b). Starting from $X_0>0$, all
$X_n$ are positive and the sequence is decreasing.

\emph{(d)} Write the equation as $n^3(X_{n-1}-X_n)=(n+1)^3(X_n-X_{n+1})+4(2n+1)^3X_n$. By (c) both
terms on the right are positive, and $(2n+1)^3>8n^3$. Hence $n^3(X_{n-1}-X_n)>32n^3X_n$, that is,
$X_{n-1}>33X_n$.
\end{proof}

\begin{lemma}[Denominators]\label{lem:lcm}
$N!=\prod_{k=1}^{N}D_{\lfloor N/k\rfloor}$ for every $N\ge1$. Consequently $D_n<3^n$ for infinitely
many $n$.
\end{lemma}
\begin{proof}
A prime $p$ divides $D_m$ exactly once for each power $p^a\le m$. So it divides the product once for
each pair $(k,a)$ with $p^a\le N/k$, that is, $\sum_a\lfloor N/p^a\rfloor$ times: as often as it
divides $N!$. Now suppose $D_m\ge3^m$ for all $m\ge M$. For $N$ large let $L=\lfloor N/M\rfloor$; the
factors with $k\le L$ have index at least $M$, so
\[
N\log N\ge\log N!\ge\log3\sum_{k\le L}\Bigl\lfloor\frac Nk\Bigr\rfloor\ge
\log3\,\Bigl(N\sum_{k\le L}\frac1k-L\Bigr)\ge\log3\cdot N\,\bigl(\log N-\log M-\tfrac1M\bigr),
\]
using $\sum_{k\le L}\frac1k\ge\log(L+1)>\log N-\log M$. Since $\log3>1$ this is false for large $N$.
\end{proof}

\begin{theorem}
$0<I_n\le2\zeta(3)\,33^{-n}$ for all $n\ge0$, and $\zeta(3)$ is irrational.
\end{theorem}
\begin{proof}
By Lemmas \ref{lem:arith}, \ref{lem:rec} and \ref{lem:bdd}, the sequence $I_n$ is bounded, starts at
$I_0=2\zeta(3)>0$ and satisfies \eqref{star}. Lemma~\ref{lem:max} gives $0<I_n<I_{n-1}/33$. If
$\zeta(3)=P/Q$ with positive integers $P,Q$, then by Lemma~\ref{lem:arith} the number
$\frac12Q\,D_n^3I_n=D_n^3a_nP-D_n^3b_nQ$ is an integer; it is positive; and for the infinitely many
$n$ of Lemma~\ref{lem:lcm} it is less than $Q\zeta(3)\,(27/33)^n$. For large such $n$ this is less
than $1$.
\end{proof}

\paragraph{Remarks.}
\begin{enumerate}
\item \emph{How sharp is $33$?} The proof gives $I_{n-1}/I_n>1+4(2+\frac1n)^3$, that is $109$,
$63.5$, $51.8$ for $n=1,2,3$. The true ratios are $116.9$, $66.8$, $54.1$, decreasing to
$17+12\sqrt2=33.97$.
\item \emph{What is not used:} integrals; the recurrence or the growth of the integers $a_n$; the
identity $a_nb_{n+1}-a_{n+1}b_n=6/(n+1)^3$ of Ap\'ery's proof; a second series; the prime number
theorem.
\item \emph{Other constants.} The same lemma, with the corresponding recurrence, gives
$0<y_n<y_{n-1}/5$ for the classical combinations of $1$ and $\log2$ (and $5>3$ suffices). For
$\zeta(2)$ the small solution of the recurrence alternates in sign; the analogous argument applied
to $(-1)^n$ times it gives a factor $11$ per step (and $11>9$ suffices). These two statements
were found by one of the agents and checked numerically for $n\le40$; they are not proved here.
\item \emph{Relation to the literature} (from a short search on 14 November 2026; passages read in
the first two items). The condition that makes Lemma~\ref{lem:max} work, namely that the middle
coefficient of Ap\'ery's recurrence exceeds the sum of the outer ones, was observed by
A.~Mingarelli (arXiv:math/0608577), who calls the recurrence \emph{disconjugate} and identifies
$a_n\zeta(3)-b_n$ as its recessive solution; he states that this is a restatement of Ap\'ery's proof
and takes the growth of $a_n$ from elsewhere. The formal proof of Mahboubi and Sibut-Pinote
(arXiv:1912.06611) contains the bound $33^n=O(a_n)$ for the integers $a_n$ and the comparison
$27<33$. K.~Rajkumar (arXiv:1212.5881) obtains the rate from the recurrence for the combinations
alone, using the Poincar\'e--Perron theorem. What we did not find is the elementary argument above,
applied directly to $I_n$: it replaces the Poincar\'e--Perron step by six lines and needs nothing
about $a_n$. The idea is therefore known; at most the packaging is new, and our search was brief.
\item \emph{Checks.} Lemma~\ref{lem:max} was proved independently by two AI agents and read line by
line by a third; identity \eqref{cert} was verified with $n$ symbolic; the conclusions $I_n>0$ and
$I_n<I_{n-1}/33$ were confirmed with exact $a_n$, $b_n$ and 150-digit arithmetic for $n\le40$.
\end{enumerate}
\end{document}
